\chapter{Quantum status}
\label{ch:quantum}

The working model of this chapter is the temporal half of \cref{ch:mass}, the Yang--Mills field $b$ of its copy symmetry, and the clock
$\theta$. Units are $\hbar=c=1$. Fields are normalised canonically, so that $m$ drops out, and $I_0=f^2$ with $f$ an energy. We expand
about the classical solution $X=b_\mu=0$, $\bar\theta=\kappa t$. The free matter spectrum and the Compton amplitude of order $g^2$ use
this background with a prescribed clock. In \cref{sec:clockamp} the clock is dynamical and $\pi=f(\theta-\bar\theta)$ is its canonically
normalised fluctuation, to which the action at $X=0$ gives the massless kinetic term; \cref{prop:ladder} quantises the homogeneous
dynamical clock at fixed rotor charge. Amplitudes are computed at tree level from the Lagrangian.

The criterion $|\mathcal M|<8\pi$, by which a growing amplitude is turned into an energy, is an order-of-magnitude criterion: it locates
where a fixed-angle tree amplitude reaches a conventional value. It is not a partial-wave unitarity bound, which would need an angular
projection in the coupled helicity channels.

The free temporal half is a quantum theory with a positive-norm Fock space. With $b$ or with the clock its tree amplitudes grow with
energy: at fixed angle they reach $8\pi$ at about $5\kappa/g$ with $b$ and at about $1.9\,(\kappa^2I_0)^{1/4}$ with the clock. A growing
tree amplitude marks a limitation of the weakly coupled model as written, which must then be read as an effective theory; it does not
settle whether a nonperturbative completion exists. In this model the clock does not remove the growth.

\section{The free theory}

\begin{proposition}[Poles and norms]
\label{prop:poles}
For each of the three theories the propagator, the inverse of the quadratic form in Fourier space, has poles only at
$\omega^2=k^2+\kappa^2$, where the quadratic form has zero eigenvalues. The residue there has rank three with all directions positive for
the temporal half, rank three with all directions negative for the spatial half, and rank six with three of each sign for the full theory.
Along a ray of fixed non-null direction, with $|p^2|\to\infty$, the propagator of the full theory falls off like $1/p^2$, while that of the
temporal half tends to a nonzero matrix of rank three and of size $1/(m\kappa^2)$. On the light cone $\omega=k$ the propagator of the full
theory is the constant $-2\eta/(m\kappa^2)$, and that of the temporal half grows like $k^2/(m\kappa^4)$. Of the three polarisations at
momentum $k$, the one with $\bB$ along $\mathbf k$ stays of size one, and the two with $\bB$ across $\mathbf k$ grow like $k/\kappa$.
\checks{qs-1, qs-2, qs-3, qs-4, qs-22}
\end{proposition}

\begin{proof}[Proof of the signs]
For $\mathbf k$ along $e_3$ the residue of the temporal half in the variable $\omega^2-|\mathbf k|^2-\kappa^2$ is
$(2/m)\sum_hv_hv_h^\top$, where $v_h$ are the three polarisations, normalised by $\eta(v_h,v_h)=1$; it is therefore positive semidefinite of
rank three. A rotation gives the other directions of $\mathbf k$, and the mirror map of \cref{prop:mirror} reverses the sign for the spatial
half. The residue of the full theory is $2\eta/m$.
\end{proof}

The temporal half therefore has three positive-norm polarisations and nothing else, and with the positive energy \cref{eq:energyB} it has a
Fock space with a ground state: the obstruction of \PI{Cor.~4.15} belongs to the full theory, where the spatial oscillators are present.
The price is the usual one. Removing the negative-norm half leaves a propagator that does not fall off along non-null rays; the full theory
has the opposite pair of properties.

\section{With the connection}

Consider two fields, one abelian $b$ and the coupling $g$; the two fields form one complex two-form of charge $g$, and
$M=\kappa$ is its mass.

\begin{theorem}[The amplitude $X+b\to X+b$]
\label{thm:compton}
At second order in $g$:
\begin{enumerate}[label=(\roman*),leftmargin=2.2em,itemsep=1pt]
\item the amplitude satisfies the Ward identity at the sampled point, to $10^{-10}$, and tends to the Thomson value
$-2g^2(\varepsilon_1\!\cdot\!\varepsilon_2^*)(\beta_2^*\!\cdot\!\beta_1)$ at low energy;
\item at $90^\circ$ in the centre-of-mass frame, with momentum $P$, the largest helicity amplitude grows as $|\mathcal M|\simeq g^2(P/\kappa)^2$,
with the fitted coefficient $1.000$ and exponent $2.000$ between $P=300\kappa$ and $1000\kappa$. Longitudinal matter states grow as well:
with both matter helicities zero and opposite helicities of $b$,
\[
\frac{|\mathcal M|}{g^2}=\frac{(v^2-1)(2v^4+7v^2-3)}{8v^2(v^2+1)},\qquad v=\frac{E+P}{\kappa},\quad E=\sqrt{\kappa^2+P^2},
\]
which grows like $(P/\kappa)^2$ with the same coefficient;
\item for a scalar of the same charge and mass, put $D=(s-M^2)(M^2-u)>0$. The two independent helicity amplitudes have the magnitudes
$2g^2(M^4-su)/D$ and $2g^2(-M^2t)/D$, which are nonnegative and add up to $2g^2$: $|\mathcal M|\le2g^2$ at all energies and angles;
\item on forty random momenta, angles and polarisations the amplitude equals that of a vector field with minimal coupling, with relative
difference below $10^{-8}$.
\end{enumerate}
\checks{qs-5, qs-6, qs-7, qs-20, qs-8, qs-21, qs-9; (ii) and (iii) numerically, the closed forms at sampled points}
\end{theorem}

The $90^\circ$ amplitude reaches $8\pi$ at $P\approx\sqrt{8\pi}\,\kappa/g\approx5\kappa/g$, a fixed-angle estimate of the scale at which
perturbation theory fails. The equality (iv) has an exact reason: $W=\star(D\wedge X)/\kappa$ obeys the minimally coupled Proca equation,
because no two covariant derivatives are commuted in the map. In the variables of \cref{ch:mass}, $W$ is $\ell$ times the magnetic
four-current, and the gyromagnetic ratio is one \cite{LeeYang62,Holstein2006}.

\begin{proposition}[What removes the growth in this amplitude]
\label{prop:moment}
Add the magnetic-moment term $i\,g\,c\,F^{\mu\nu}\bar W_\mu W_\nu$ to the vector form, of mass $M=\kappa$.
\begin{enumerate}[label=(\roman*),leftmargin=2.2em,itemsep=1pt]
\item At $90^\circ$: for $c=0$ the largest amplitude grows with exponent two; for $c=+1$, gyromagnetic ratio two, it is bounded at this
angle and tends to $4g^2$; the longitudinal channel of \cref{thm:compton}(ii), with opposite helicities of $b$, vanishes at this angle, and
the longitudinal channels with equal helicities of $b$ have the magnitude $2g^2P/E$, which tends to $2g^2$; for $c=-1$ the largest amplitude grows
with the same exponent as for $c=0$ and four times its coefficient.
\item For $c=+1$ the growth cancels at each fixed angle $0<\vartheta<\pi$, with a limit that depends on the angle. The amplitude with all
helicities $+1$, for which Bjerrum-Bohr, Holstein, Plant\'e and Vanhove give the spin-one helicity amplitudes
\cite[Appendix~A]{BjerrumBohr2015}, has the magnitude $2[(s-M^2)^2+M^2t]^2/\bigl((s-M^2)^3(M^2-u)\bigr)\,g^2$, and at fixed $\vartheta$ it tends to
$4g^2/(1+\cos\vartheta)$: $4g^2$ at $90^\circ$, $8g^2$ at $120^\circ$. At $s=100M^2$ and $120^\circ$ it is already $7.651\,g^2$, so no bound
$4g^2$ holds at all angles.
\item The backward direction is exceptional. At $\vartheta=\pi$, where $t=-(s-M^2)^2/s$ and $u=M^4/s$, the same expression is
$2(s-M^2)^2g^2/(sM^2)$, which grows like $2s/M^2$: $800$, $8\times10^4$ and $8\times10^6$ times $g^2$ at $P/M=10$, $100$ and $1000$. The
limit at fixed angle cannot be taken at $\vartheta=\pi$; the growth is a backward peak of angular width of order $M/\sqrt s$.
\end{enumerate}
\checks{qs-10, qs-17, qs-23, qs-19, numerically, (ii) and (iii) against the closed form}
\end{proposition}

The value two is the one that tree unitarity at fixed angle requires \cite{FerraraPorratiTelegdi92} and the one that a Yang--Mills
embedding supplies (\cref{prop:whatadds}).

\begin{proposition}[Several fields]
\label{prop:stueck}
For $\so3$ and a copy-vector-valued one-form $A$, $D\wedge(D\wedge A)=g\,F_b\wedge A$. So the substitution $X\to X+D\wedge A/\kappa$, which
restores the two-form gauge symmetry, changes $H=D\wedge X$ by $(g/\kappa)F_b\wedge A$ and exposes the interactions
$(g/\kappa)(D\wedge X)\cdot(F_b\wedge A)$ and $(g/\kappa)^2(F_b\wedge A)^2$, whose scale is again $\kappa/g$. With $Y=X+D\wedge A/\kappa$
and the connection held fixed, the transformations $\delta X=D\wedge\Lambda$, $\delta A=-\kappa\Lambda$ leave $Y$ unchanged: this is the
redundancy that the substitution introduces.
\checks{qs-12, for the identity; no amplitude was computed for three fields}
\end{proposition}

\section{With the clock}
\label{sec:clockamp}

Here the clock is dynamical: $\theta=\kappa t+\pi/f$, so that $(\partial\theta)^2=\kappa^2+(2\kappa/f)\,\partial_t\pi+(\partial\pi)^2/f^2$,
and $\pi$ is a massless field.

\begin{theorem}[The amplitude $X+\pi\to X+\pi$]
\label{thm:clockamp}
At $90^\circ$ in the frame of the clock, label the amplitudes by the helicities $(h_{\rm in},h_{\rm out})$ of the incoming and the outgoing
matter state. Between transverse states, $h_{\rm in},h_{\rm out}=\pm1$, the amplitudes grow as $|\mathcal M|\simeq2P^4/(\kappa^2f^2)$, with
the fitted coefficient $2.000$ and exponent $4.000$. Between longitudinal states, $(0,0)$, the amplitude grows as $|\mathcal M|\simeq2P^2/f^2$.
The mixed channels grow more slowly, linearly in $P$ at $\kappa=f$.
\checks{qs-11, numerically}
\end{theorem}

The transverse amplitude reaches $8\pi$ at $P\approx(4\pi)^{1/4}\sqrt{\kappa f}\approx1.9\,(\kappa^2I_0)^{1/4}$, again a fixed-angle
estimate. Since $I_0\kappa^2/2$ is the energy density of the clock, this is about twice the fourth root of that energy density.
Dimensionally, it is the scale at which the fluctuation $\partial_t\pi/f$ of the rate, probed at the scale $1/P$, would be comparable to the
rate itself, which suggests that the degenerate case of \cref{prop:clockhalf} becomes relevant there; that fluctuation, which would need a
state, a smearing scale and an observable, is not computed. A window $\kappa\ll P\ll1.9\sqrt{\kappa f}$ exists only if $I_0\gg\kappa^2$.

\begin{proposition}[The uniform sector as quantum mechanics]
\label{prop:ladder}
For fields that do not depend on $\mathbf x$, with the volume absorbed into the mechanical parameters $m$ and $I_0$, the reduced Hamiltonian
is $H=|\bp|^2/2m+h^2/\bigl(2(I_0-m|\bB|^2)\bigr)$, with $h=\hbar N$ conserved, on the domain $|\bB|<\sqrt{I_0/m}$ where the inertia is
positive. Quantise it canonically, with the flat measure in $\bB$; for a compact rotor angle $N$ is an integer. Take $N>0$ (for $N<0$
replace $N$ by $|N|$; $N=0$ lies outside the expansion), Dirichlet conditions at $|\bB|=\sqrt{I_0/m}$, and the s-wave sector. Then $H$ is
bounded below, and with $\kappa_N=\hbar N/I_0$, for each fixed radial level as $N\to\infty$:
\begin{enumerate}[label=(\roman*),leftmargin=2.2em,itemsep=1pt]
\item the levels above the rotor energy are $\hbar\kappa_N(1.5,3.5,5.5,\dots)(1+O(1/N))$;
\item the ground state has the asymptotic expansion
\[
E_0(N)=\frac{\hbar^2}{I_0}\Bigl(\frac{N^2}2+\frac{3N}2+\frac{15}8+\frac{45}{32N}+O(N^{-2})\Bigr),
\]
and at $N=1600$ the shift of the lowest level times $N$ is $1.876$, as $15/8+45/(32N)$ gives;
\item the ladder of \PI{Prop.~4.52}(i) holds with a zero-point term, up to a correction of order $\hbar^2/(I_0N^2)$:
$E_0(N+1)-E_0(N)=\hbar(\kappa_N+\kappa_{N+1})/2+\tfrac32\hbar(\kappa_{N+1}-\kappa_N)+O(\hbar^2/(I_0N^2))$.
\end{enumerate}
\checks{qs-15, qs-16, qs-18; (ii) exactly by perturbation theory, and numerically}
\end{proposition}

In this proposition the mechanical inertia $I_0$ has absorbed the spatial volume $V$; write it $I_{\rm mech}=f^2V$ for the comparison. The
mechanical expansion is controlled by $1/N=\hbar/(I_{\rm mech}\kappa)$, which is $1/(f^2V\kappa)$ for $\hbar=c=1$, the quantum form of the
multiscale criterion of \PI{Cor.~4.30}. The field theory has the dimensionless ratio $\kappa^2/f^2$ at momenta of order $\kappa$, which
corresponds to this mechanical parameter in a volume of order $\kappa^{-3}$; at large momentum the transverse clock amplitude grows as
$P^4/(\kappa^2f^2)$ (\cref{thm:clockamp}).

\begin{proposition}[What protects the relation between mass and rate]
\label{prop:shifts}
$D_\mu(e^{-\theta G}c)=0$ for every constant $c\in\R^{3,3}$, so the full clock-coupled theory of \cref{thm:clock} is invariant under the six
shifts $X\to X+e^{-\theta G}c$, the frame translations of the group of \PI{Prop.~4.24}. They forbid a mass term that does not come from the
clock. The form $\eta(X,X)$ is not invariant under them, so the mass term of the temporal half breaks them.
\checks{qs-13, qs-14}
\end{proposition}

\section{Bearing on statements of Part~I}

A quantum theory of the momentum field exists for the temporal half (\cref{prop:poles,prop:ladder}); \PI{Cor.~4.15 and Section~4.5.8}
concern the full theory and stand. The zero-point cancellation of \PI{Thm~9.12} needs three positive and three negative oscillators at each
frequency, and the temporal half keeps only the positive ones, so its zero-point energy does not cancel. The rate ladder of
\PI{Prop.~4.52} survives in the mechanical model (\cref{prop:ladder}); in the field theory the clock background is a state of constant
charge density of the shift symmetry of $\theta$, and the rate is not discrete, the spacing $\hbar/(I_0\times\text{volume})$ going to zero.
The window $\kappa<P<5\kappa/g$ needs the coupling of $b$ to be small at the scale $\kappa$. If $\hbar\kappa=\Lambda_b$, the scale at which
a non-abelian $b$ becomes strong, as \PI{Conjecture~9.17} proposes, a weak-coupling treatment at the scale $\kappa$ is not available, and
the tree-level analysis of this chapter does not apply there; how strong the coupling is, and whether the massive quanta are then
composites of $b$, would need the running and the matter content, which are not specified here. Finally, once $b$ or the clock is coupled,
the remaining symmetries allow an electric-current kinetic term $J_e\!\cdot\!J_e$ of the spatial half. Its coefficient and sign, and the
existence and nature of any additional pole within the range of the effective theory, require a matching or a loop calculation, which is
not done: with one sign the term adds a pole with a negative kinetic term, with the other a tachyonic mode, and without it $\bE$ stays
auxiliary.

\section{What a completion needs}
\label{sec:needs}

In the model of this chapter the clock does not remove the growth. The amplitude with the clock grows in the channels computed
(\cref{thm:clockamp}), and in the minimal model about $X=0$ there is no tree-level vertex coupling $\pi$ to two quanta of $b$, so the clock
supplies no cancellation of the order-$g^2$ amplitude of \cref{thm:compton}. The growing channels there include transverse and longitudinal
matter states (\cref{thm:compton}(ii)). These are statements about this model, not a general exclusion of scalar interactions. For one
neutral field a vector with a coupling of the form $X\wedge F$ gives the mass \cite{CremmerScherk74,AllenBowickLahiri91}, with a constant
mass and no clock.

For charged massive spin-one particles, Llewellyn Smith gave the first argument \cite{LlewellynSmith73}, and Cornwall, Levin and
Tiktopoulos the general result \cite{CornwallLevinTiktopoulos74}. They impose the high-energy bound on all tree amplitudes of a specified
class of Lagrangians of scalars, spinors and vectors, with interactions of dimension at most four, at fixed nonexceptional kinematics. Under
their hypotheses the theory is equivalent, up to field redefinitions and with their exception for masses associated with invariant abelian
subgroups, to a spontaneously broken gauge theory, in which the massive charged vectors and $b$ are components of one connection. A bound
in one helicity channel or at one angle, as computed here, does not by itself establish these hypotheses. For $n$ massive fields in the vector of $\so n$ a compact embedding is $\so{n+1}\supset\so n$: the adjoint of $\so{n+1}$
decomposes into the unbroken $\so n$ and one $n$-vector, which is why it works. Other embeddings, with further generators, spectator factors
or scalar sectors, are not excluded by this argument; we choose the minimal orthogonal one, $\so{n+1}$ broken to $\so n$. For $n=2$ it is
$\so3$ broken to $\so2$, the pattern that \PI{Remark~9.63}(c) finds for the temporal rotations with the clock as the Higgs vector. In such a
completion the mass is $g$ times the length of a vector field, and not the rate of $\theta$.
